\section{Related Work}
\label{sec:related}

\paragraph{Algorithmic collusion.} \citet{calvano2020} show that Q-learning
agents in a repeated Bertrand game learn supracompetitive prices sustained by
punishment, documented through impulse responses to a forced deviation;
\citet{klein2021} obtains similar results under sequential pricing. A second
wave questions how these outcomes arise. \citet{asker2022,asker2024} show that
the outcome depends on algorithm design: asynchronous updating (learning only
from the action taken) yields supracompetitive prices, while synchronous,
counterfactual updating yields competitive ones. \citet{abada2023} show that
algorithms quickly reach seemingly collusive outcomes and that a regulator can
restore competition by enforcing decentralised learning or intervening in the
learning process.
\citet{banchio2023} show that collusion can arise from spontaneous coupling
of the learning dynamics rather than from punishment strategies.
\citet{denboer2026} argue that the collusive equilibria of Q-learning are
reached only on time scales irrelevant in practice and require implausibly
synchronized algorithms. \citet{calvano2023} discuss when algorithmic
collusion is genuine or spurious. Deep RL agents \citep{deng2024} and large
language models \citep{fish2024} can also price above competitive levels.
We contribute a placebo-based diagnostic (the myopic learner) and a
detectability analysis to this debate, in a financial market with
adverse selection.

\paragraph{Algorithms in financial markets.} \citet{colliard2026} study
Q-learning market makers and find mark-ups over the competitive price driven
by limited experimentation and noisy feedback, the same mechanism we find for
informed traders. \citet{cartea2022} show that the tick size bounds the
excess rents of learning liquidity providers, and \citet{cont2024} study
tacit collusion among market-making algorithms in dealer markets.
\citet{dou2025} extend \citet{kyle1985} to several informed speculators,
information-insensitive investors, and a market maker who trades off pricing
error against inventory. They show that Q-learning speculators sustain
collusive profits and reduce price informativeness. Their theory shows that
price-trigger equilibria exist only when information-insensitive investors
are numerous and noise trading is low; in simulations they separate
price-trigger collusion from over-pruning with noise-shock and deviation
impulse responses, finding the former in that region and the latter
elsewhere, and find that removing the lagged price from the state
eliminates collusion in the price-trigger region. An
independent reimplementation \citep{esquinas2026} reproduces both regimes and
finds that the level of collusion depends on the update protocol. Our paper
adds the rival-deviation test, the myopic placebo and the learning-rate
sensitivity analysis, quantifies the monitoring problem behind their regime
split, and provides a validated open testbed.

\paragraph{Estimation bias in Q-learning.} The max operator makes Q-learning
over-estimate action values under noise \citep{thrun1993,vanhasselt2010}. The
bias relevant here runs the other way for rarely chosen actions: with a
constant step size and $\varepsilon$-greedy exploration, an action whose
estimate falls below the greedy action by chance is rarely updated again, so
pessimistic errors persist while optimistic ones are corrected. Actions with
noisier payoffs are more exposed. Decreasing step sizes that satisfy the
\citet{robbins1951} conditions remove the persistent component
asymptotically \citep{evendar2003}.
